\section*{الفصل \ref{ch:b3:green-industrial} --- الكيمياء الخضراء والعمليات الصناعية}

\begin{solution}{exo:b3:green-industrial:1}
ديلز--ألدر: 100\,\% (إضافة). الأسترة: $88.0/(60.0 + 46.0) = 83$\,\%. فيتيغ:
$104.0/(106.0 + 276.0) = 27$\,\%، إذ يحمل أكسيد ثلاثي فينيل الفوسفين معظم الكتلة.
\end{solution}

\begin{solution}{exo:b3:green-industrial:2}
$\mathrm{PMI} = 120/8.0 = 15$؛ $E = \mathrm{PMI} - 1 = 14$.
\end{solution}

\begin{solution}{exo:b3:green-industrial:3}
غسلة واحدة لحمولة قياسية عند درجة حرارة معينة ونتيجة تنظيف معينة. فالمنظفات
تُجرَّع بمقادير مختلفة: المسحوق المركّز يحتاج إلى كمية أقل في الغسلة، فلا تقدّم الكتل
المتساوية الخدمة نفسها.
\end{solution}

\begin{solution}{exo:b3:green-industrial:4}
2-ميثيل رباعي هيدروفوران أو أسيتات الإيثيل مكان ثنائي كلوروميثان؛ والهبتان (مع أسيتات
الإيثيل) مكان الهكسان؛ ومكان ثنائي ميثيل فورماميد: أسيتات الإيثيل أو 2-ميثيل رباعي هيدروفوران أو
كربونات ثنائي الميثيل، أو اقتران في الماء بمواد فعالة سطحيًا، بحسب
الكواشف.
\end{solution}

\begin{solution}{exo:b3:green-industrial:5}
1.000 mol من الحمض؛ و0.800 mol من الإستر (\qty{88.0}{g/mol}): المردود 80\,\%. AE يساوي 83.0\,\%. وSF $= 152/106 = 1.43$.
$\mathrm{RME} = 70.4/152 = 0.463$؛ و$\mathrm{AE} \times \text{المردود}/\mathrm{SF} = 0.830 \times 0.80/1.43 = 0.463$.
\end{solution}

\begin{solution}{exo:b3:green-industrial:6}
الكلوروهيدرين: $44.0/(28.0 + 71.0 + 74.1) = 25$\,\%؛ الأكسدة المباشرة: 100\,\%. مول واحد من $\ce{CaCl2}$ لكل مول من
أكسيد الإيثيلين: $\qty{1.0e6}{g}/\qty{44.0}{g/mol} \times \qty{111.1}{g/mol} = \qty{2.5}{t}$ لكل طن.
\end{solution}

\begin{solution}{exo:b3:green-industrial:7}
$3/8 \times \qty{5.88e4}{mol} \times \qty{44.0}{g/mol} = \qty{0.97}{t}$ من $\ce{CO2}$ لكل طن من الأمونيا. وللهيدروجين، $\ce{CH4 + 2H2O -> CO2 + 4H2}$:
$1/4$ مول لكل مول، و$\qty{5.0e5}{mol}/4 \times \qty{44.0}{g/mol} = \qty{5.5}{t}$ من $\ce{CO2}$ لكل طن من الهيدروجين.
\end{solution}

\begin{solution}{exo:b3:green-industrial:8}
% ledger: dfg:H2O_l, dfh:H2O-l
الكيلوغرام الواحد يساوي \qty{500}{mol}. من $\Delta_rG^\circ = \qty{237.1}{kJ/mol}$: $\qty{1.19e5}{kJ} = \qty{32.9}{kWh}$. ومن $\Delta_rH^\circ = \qty{285.8}{kJ/mol}$: \qty{39.7}{kWh}
(والفرق حرارة يمكن أن يقدّمها الوسط المحيط).
\end{solution}

\begin{solution}{exo:b3:green-industrial:9}
$\qty{1.0e6}{g}/\qty{146.0}{g/mol} = \qty{6.85e3}{mol}$ من $\ce{N2O}$، أي $\qty{0.30}{t}$؛ و$\times 273 = \qty{82}{t}$ من مكافئ $\ce{CO2}$ لكل طن من حمض الأديبيك.
\end{solution}

\begin{solution}{exo:b3:green-industrial:10}
تحتسب \omterm{def:b3:green-industrial:rme}{كفاءة كتلة التفاعل} لخطوة ما الكاشفَ الفائض مستهلَكًا؛ فإن استُرجع الفائض
وأُعيد إلى التغذية، لم يُستهلك إلا الجزء الضائع فعلًا، وتكون كفاءة الكتلة
الإجمالية أعلى من جداء قيم الخطوات. وفي \omterm{def:b3:green-industrial:e-factor}{العامل E} لا يُعد
التيار المعاد تدويره نفايات، شرط أن تكون إعادة التدوير داخل الحدود؛
وعندئذ يجب احتساب طاقتها وفواقدها.
\end{solution}

\begin{solution}{exo:b3:green-industrial:11}
قبل التحسين: $20 \times 0.80 = \qty{16}{kg}$ من المذيب، يضيع منها 10\,\%: أي \qty{1.6}{kg} من النفايات. وبعده: \qty{4.0}{kg}، يضيع منها \qty{0.40}{kg}. فينخفض
\omterm{def:b3:green-industrial:e-factor}{العامل E} الكامل بمقدار 1.2 kg لكل كيلوغرام من الناتج (وتنخفض طاقة التقطير
بثلاثة أرباع).
\end{solution}

\begin{solution}{exo:b3:green-industrial:12}
لكل \omterm{def:b3:green-industrial:lca}{وحدة وظيفية} وبالحدود نفسها، يعرض كل فئة أثر
على حدة (المناخ، واستعمال الأرض، واستعمال الماء\dots) بدلًا من درجة واحدة. ويتطلب
القرار الأوزانَ المعطاة لهذه الفئات، وندرة الماء
والأرض محليًا، وعدم يقين البيانات، ومعرفة ما إذا كانت المادة الأولية الحيوية
تنافس الغذاء.
\end{solution}

\begin{solution}{pb:b3:green-industrial:1}
\textbf{1.} $\ce{C10H14 + C4H6O3 -> C12H16O + C2H4O2}$ (HF)؛ $\ce{C12H16O + H2 -> C12H18O}$ (نيكل راني)؛
$\ce{C12H18O + CO -> C13H18O2}$ (البلاديوم).

\textbf{2.} $\ce{C10H14 + C4H6O3 + H2 + CO -> C13H18O2 + C2H4O2}$.

\textbf{3.} $134.0 + 102.0 + 2.0 + 28.0 = \qty{266.0}{g/mol}$ من المتفاعلات؛ والإيبوبروفين \qty{206.0}{g/mol}: AE $= 206.0/266.0 = 77$\,\%.

\textbf{4.} الإيزوبوتيل بنزين، وأنهيدريد الأسيتيك، وكلورو أسيتات الإيثيل، وإيثوكسيد الصوديوم،
والحمض المائي ($\ce{H3O+}$)، والهيدروكسيل أمين والماء.

\textbf{5.} $134.0 + 102.0 + 122.5 + 68.0 + 19.0 + 33.0 + 2 \times 18.0 = \qty{514.5}{g/mol}$: AE $= 206.0/514.5 = 40$\,\%.

\textbf{6.} حمض الأسيتيك، وكلوريد الصوديوم، والإيثانول، وثنائي أكسيد الكربون، والماء، والأمونيا (في صورة
ملح أمونيوم)، وأملاح الألومنيوم الآتية من كلوريد الألومنيوم المستعمل بنسبة ستوكيومترية.

\textbf{7.} $0.71 + 0.55 + 0.010 + 0.14 + 0.05 = \qty{1.46}{t}$.

\textbf{8.} كثافة الكتلة 1.46؛ $E = 0.46$.

\textbf{9.} $0.46 - 0.31 = 0.15$.

\textbf{10.} كثافة الكتلة PMI $= 2.0 + 1.5 = 3.5$؛ $E = 2.5$.

\textbf{11.} المردودات دون 100\,\%، والكواشف الفائضة، وفواقد المذيب والحفاز، 
ومواد المعالجة كلها تضيف نفايات يتجاهلها \omterm{def:b3:green-industrial:atom-economy}{الاقتصاد الذري}.

\textbf{12.} منع النفايات، وتعظيم \omterm{def:b3:green-industrial:atom-economy}{الاقتصاد الذري}، واستعمال الحفازات بدلًا من
الكواشف المستعملة بنسب ستوكيومترية، وتجنب المشتقات والخطوات غير الضرورية، وتوفير الطاقة.

\textbf{13.} كلوريد الألومنيوم المستعمل بنسبة ستوكيومترية، الذي يتحلمأ إلى نفايات ألومنيوم عند
المعالجة؛ أما فلوريد الهيدروجين فهو حفاز ومذيب معًا، ويُسترجع ويُعاد استعماله.

\textbf{14.} نيكل راني (حفاز هدرجة غير متجانس).

\textbf{15.} معقد فوسفيني للبلاديوم: كيمياء \omterm{def:b3:homogeneous-catalysis:carbonylation}{الكربونلة} في
\cref{ch:b3:homogeneous-catalysis}.

\textbf{16.} ينقل أنهيدريد الأسيتيك مجموعة أسيتيل واحدة؛ ويغادر النصف الآخر من
الجزيء في صورة حمض الأسيتيك.

\textbf{17.} فلوريد الهيدروجين شديد السمية وأكّال، وأول أكسيد الكربون سام 
وقابل للاشتعال: معدات مغلقة من مواد مقاومة، وأجهزة كشف، ومشغّلون
مدرَّبون.

\textbf{18.} كل خطوة تضيف تسخينًا وتبريدًا وعمليات فصل واسترجاعًا للمذيب؛ والخطوات
الثلاث تحتاج إلى عدد أقل منها كلها.

\textbf{19.} $2.5 \times 3000 = \qty{7500}{t}$ في السنة.

\textbf{20.} $0.15 \times 3000 = \qty{450}{t}$ في السنة.

\textbf{21.} نحو \qty{7000}{t} من النفايات المتجنَّبة كل سنة.

\textbf{22.} لا: فطبيعة النفايات مهمة (فطن من الماء المالح ليس كطن من
مركب سام ثابت)، وكذلك الطاقة والانبعاثات ومخاطر
الكواشف.

\textbf{23.} آثار صنع الكواشف والطاقة في المراحل السابقة، والانبعاثات إلى
الهواء والماء، واستعمال المنتج والتخلص منه، لكل \omterm{def:b3:green-industrial:lca}{وحدة وظيفية}، في
عدة فئات أثر.

\textbf{24.} للطريق الحفزي ذي الخطوات الثلاث \omterm{def:b3:green-industrial:atom-economy}{اقتصاد ذري} نحو 77\,\% (مقابل 40\,\% للطريق
ذي الخطوات الست).
\end{solution}
